\s{
    \section{Classification of electrical machines}
    Electrical machines is another term for electromagnetic energy converters and describes, on the one hand, transformers as stationary electrical machines and, on the other hand, electric motors and generators.
    Electric motors power many everyday devices such as household appliances and electric vehicles, while generators and transformers ensure the generation and distribution of electrical energy.
 
    Dieses Kapitel bietet eine Einführung in die grundlegenden Konzepte und Funktionsweisen elektrischer Maschinen. Zunächst werden die verschiedenen Typen elektrischer Maschinen vorgestellt, einschließlich ihrer Prinzipien der Energieumwandlung und ihrer Hauptkomponenten. Es folgt eine detaillierte Betrachtung der Funktionsweise von Elektromotoren und Generatoren, der Relevanz von Transformatoren für die Energieübertragung sowie der Einfluss neuer Technologien auf die Effizienz und Leistung elektrischer Maschinen. 

    \begin{Learning objectives}{Elektrische Maschinen}
        The students
        \begin{itemize}
            \item are familiar with the various types of electrical machines.
            \item can solve basic problems in the field of electrical machines.
        \end{itemize}
    \end{Learning objectives}
}
\b{
    \begin{frame}{Learning objectives}
        \begin{Learning objectives}{Electrical machinery}
            The students
            \begin{itemize}
                \item are familiar with the various types of electrical machines.
                \item can solve basic problems in the field of electrical machines.
            \end{itemize}
        \end{Learning objectives}
    \end{frame}
}
\s{
    Electrical machines are energy converters. A distinction is primarily made between transformers and motors or generators. Transformers convert electrical energy into another form of electrical energy by transferring one voltage level to another.
    Motors and generators, on the other hand, convert electrical energy into mechanical energy and vice versa. Motors convert electrical energy into mechanical energy, while generators convert mechanical energy into electrical energy. Since the conversion 
    process is reversible, motors can also be used as generators and vice versa. Motors and generators are further 
    classified according to the type of current they operate on: direct current machines and alternating current machines. Alternating current machines are further divided into synchronous and asynchronous machines.
     
    The functioning of all electrical machines is based on the electromagnetic principles of induction and Lorentz force discussed in Module 6. 
    While the transformer is based purely on the principle of induction, the motor or generator uses different electromagnetic principles depending on the conversion process. When functioning as a generator, it relies on induction to generate electrical energy 
    through changing magnetic fields. It utilises the principle of Lorentz force when functioning as a motor.
 }


